
\documentclass{article}
\usepackage{colortbl}
\usepackage{makecell}
\usepackage{multirow}
\usepackage{supertabular}

\begin{document}

\newcounter{utterance}

\centering \large Interaction Transcript for game `hot\_air\_balloon', experiment `air\_balloon\_survival\_cs\_complexity\_hard', episode 0 with gpt{-}5.4{-}azure.
\vspace{24pt}

{ \footnotesize  \setcounter{utterance}{1}
\setlength{\tabcolsep}{0pt}
\begin{supertabular}{c@{$\;$}|p{.15\linewidth}@{}p{.15\linewidth}p{.15\linewidth}p{.15\linewidth}p{.15\linewidth}p{.15\linewidth}}
    \# & $\;$A & \multicolumn{4}{c}{Game Master} & $\;\:$B\\
    \hline

    \theutterance \stepcounter{utterance}  
    & & \multicolumn{4}{p{0.6\linewidth}}{
        \cellcolor[rgb]{0.9,0.9,0.9}{
            \makecell[{{p{\linewidth}}}]{
                \texttt{\tiny{[P1$\langle$GM]}}
                \texttt{Účastníte se hry založené na společném vyjednávání.} \\
\\ 
\texttt{Společně s dalším účastníkem se musíte dohodnout na jedné sadě předmětů, které budou ponechány. Každý z vás má vlastní představu o tom, jak moc je pro vás každý předmět důležitý (důležitost). Nevíte, jak si předměty cení druhý účastník. Navíc máte dáno, jaké úsilí každý předmět vyžaduje.} \\
\texttt{Na sadě se můžete dohodnout pouze tehdy, pokud celkové úsilí vybraných předmětů nepřekročí sdílený limit:} \\
\\ 
\texttt{LIMIT = 4734} \\
\\ 
\texttt{Zde jsou uvedeny hodnoty úsilí jednotlivých předmětů:} \\
\\ 
\texttt{Úsilí předmětu = \{"C08": 192, "C55": 325, "A43": 812, "A67": 868, "B42": 593, "A62": 915, "B19": 944, "C43": 311, "C73": 252, "C10": 343, "B36": 104, "A18": 558, "C61": 627, "C67": 593, "C42": 827, "C70": 611, "B40": 95, "C39": 251, "A21": 226, "C16": 21\}} \\
\\ 
\texttt{Zde je váš osobní pohled na důležitost každého předmětu:} \\
\\ 
\texttt{Hodnoty důležitosti předmětu = \{"C08": 286, "C55": 419, "A43": 906, "A67": 962, "B42": 687, "A62": 1009, "B19": 1038, "C43": 405, "C73": 346, "C10": 437, "B36": 198, "A18": 652, "C61": 721, "C67": 687, "C42": 921, "C70": 705, "B40": 189, "C39": 345, "A21": 320, "C16": 115\}} \\
\\ 
\texttt{Cíl:} \\
\\ 
\texttt{Vaším cílem je vyjednat sdílenou sadu předmětů, která vám přinese co největší prospěch (tj. maximalizuje celkovou důležitost pro VÁS), a zároveň zůstat v rámci LIMITU. Nejste povinni v každé zprávě podat NÁVRH {-} můžete také jednoduše vyjednávat. Všechny taktiky jsou povoleny!} \\
\\ 
\texttt{Protokol interakce:} \\
\\ 
\texttt{Ve zprávě smíte používat pouze následující strukturované formáty:} \\
\\ 
\texttt{NÁVRH: \{'A', 'B', 'C', …\}} \\
\texttt{Navrhněte ponechat přesně tyto předměty.} \\
\texttt{ODMÍTÁM: \{'A', 'B', 'C', …\}} \\
\texttt{Výslovně odmítněte návrh protivníka.} \\
\texttt{ARGUMENT: \{'...'\}} \\
\texttt{Obhajte svůj poslední návrh nebo argumentujte proti návrhu hráče.} \\
\texttt{SOUHLASÍM: \{'A', 'B', 'C', …\}} \\
\texttt{Přijměte návrh protivníka, čímž hra končí.} \\
\\ 
\\ 
\texttt{Pravidla:} \\
\\ 
\texttt{Můžete SOUHLASÍM pouze s návrhem, který druhá strana zaznamenala prostřednictvím NÁVRH.} \\
\texttt{Můžete ODMÍTÁM pouze návrh, který druhá strana zaznamenala prostřednictvím NÁVRH.} \\
\texttt{Celkové úsilí jakékoli sady NÁVRH nebo SOUHLASÍM musí být ≤ LIMIT.} \\
\texttt{NEZVEŘEJŇUJTE své skryté skóre důležitosti.} \\
\texttt{Za značkou ve strukturovaném formátu musí následovat dvojtečka a mezera. Argumentem musí být pythonová množina obsahující 0 nebo více řetězců.} \\
\texttt{Musí tedy mít tvar "TAG: \{...\}"} \\
\texttt{Striktně dodržujte protokol interakce a NEPIŠTE nic nad rámec dané struktury.} \\
\texttt{Hra končí, když jedna strana dá SOUHLASÍM k NÁVRH učiněnému druhým hráčem.} \\
\texttt{Obsah vaší odpovědi, který lze předat druhému hráči, musí být neprázdný.} \\
\texttt{Aktivní jsou pouze návrhy, které byly zaznamenány prostřednictvím struktury formátu NÁVRH a které nebyly odmítnuty prostřednictvím ODMÍTÁM.} \\
\\ 
            }
        }
    }
    & \\ \\

    \theutterance \stepcounter{utterance}  
    & \multicolumn{4}{p{0.6\linewidth}}{
        \cellcolor[rgb]{0.9,0.9,0.9}{
            \makecell[{{p{\linewidth}}}]{
                \texttt{\tiny{[P1$\rangle$GM]}}
                \texttt{NÁVRH: \{'C08', 'C55', 'A67', 'B42', 'C43', 'C73', 'C10', 'B36', 'A18', 'C67', 'B40', 'C39', 'A21', 'C16'\}} \\
            }
        }
    }
    & & \\ \\

    \theutterance \stepcounter{utterance}  
    & & \multicolumn{4}{p{0.6\linewidth}}{
        \cellcolor[rgb]{0.9,0.9,0.9}{
            \makecell[{{p{\linewidth}}}]{
                \texttt{\tiny{[GM$\rangle$P2]}}
                \texttt{Účastníte se hry založené na společném vyjednávání.} \\
\\ 
\texttt{Společně s dalším účastníkem se musíte dohodnout na jedné sadě předmětů, které budou ponechány. Každý z vás má vlastní představu o tom, jak moc je pro vás každý předmět důležitý (důležitost). Nevíte, jak si předměty cení druhý účastník. Navíc máte dáno, jaké úsilí každý předmět vyžaduje.} \\
\texttt{Na sadě se můžete dohodnout pouze tehdy, pokud celkové úsilí vybraných předmětů nepřekročí sdílený limit:} \\
\\ 
\texttt{LIMIT = 4734} \\
\\ 
\texttt{Zde jsou uvedeny hodnoty úsilí jednotlivých předmětů:} \\
\\ 
\texttt{Úsilí předmětu = \{"C08": 192, "C55": 325, "A43": 812, "A67": 868, "B42": 593, "A62": 915, "B19": 944, "C43": 311, "C73": 252, "C10": 343, "B36": 104, "A18": 558, "C61": 627, "C67": 593, "C42": 827, "C70": 611, "B40": 95, "C39": 251, "A21": 226, "C16": 21\}} \\
\\ 
\texttt{Zde je váš osobní pohled na důležitost jednotlivých předmětů:} \\
\\ 
\texttt{Hodnoty důležitosti předmětu = \{"C08": 286, "C55": 419, "A43": 906, "A67": 962, "B42": 687, "A62": 1009, "B19": 1038, "C43": 405, "C73": 346, "C10": 437, "B36": 198, "A18": 652, "C61": 721, "C67": 687, "C42": 921, "C70": 705, "B40": 189, "C39": 345, "A21": 320, "C16": 115\}} \\
\\ 
\texttt{Cíl:} \\
\\ 
\texttt{Vaším cílem je vyjednat sdílenou sadu předmětů, která vám přinese co největší prospěch (tj. maximalizuje celkovou důležitost pro VÁS), a zároveň zůstat v rámci LIMITU. Nejste povinni v každé zprávě podat NÁVRH {-} můžete také jednoduše vyjednávat. Všechny taktiky jsou povoleny!} \\
\\ 
\texttt{Protokol interakce:} \\
\\ 
\texttt{Ve zprávě smíte používat pouze následující strukturované formáty:} \\
\\ 
\texttt{NÁVRH: \{'A', 'B', 'C', …\}} \\
\texttt{Navrhněte ponechat přesně tyto předměty.} \\
\texttt{ODMÍTÁM: \{'A', 'B', 'C', …\}} \\
\texttt{Výslovně odmítněte návrh protivníka.} \\
\texttt{ARGUMENT: \{'...'\}} \\
\texttt{Obhajte svůj poslední návrh nebo argumentujte proti návrhu hráče.} \\
\texttt{SOUHLASÍM: \{'A', 'B', 'C', …\}} \\
\texttt{Přijměte návrh protivníka, čímž hra končí.} \\
\\ 
\\ 
\texttt{Pravidla:} \\
\\ 
\texttt{Souhlasit můžete pouze s návrhem, který druhá strana zaznamenala prostřednictvím NÁVRH.} \\
\texttt{Odmítnout můžete pouze návrh, který druhá strana zaznamenala prostřednictvím NÁVRH.} \\
\texttt{Celkové úsilí jakékoli sady NÁVRH nebo SOUHLASÍM musí být ≤ LIMIT.} \\
\texttt{NEZVEŘEJŇUJTE své skryté skóre důležitosti.} \\
\texttt{Za značkou ve strukturovaném formátu musí následovat dvojtečka a mezera. Argumentem musí být pythonová množina obsahující 0 nebo více řetězců.} \\
\texttt{Musí tedy mít tvar "TAG: \{...\}"} \\
\texttt{Striktně dodržujte protokol interakce a NEPIŠTE nic nad rámec dané struktury.} \\
\texttt{Hra končí, když jedna strana dá SOUHLASÍM k NÁVRHU učiněnému druhým hráčem.} \\
\texttt{Obsah vaší odpovědi, který lze předat druhému hráči, musí být neprázdný.} \\
\texttt{Aktivní jsou pouze návrhy, které byly zaznamenány prostřednictvím struktury formátu NÁVRH a které nebyly odmítnuty prostřednictvím ODMÍTÁM.} \\
\\ 
\\ 
\\ 
\texttt{Nyní obdržíte první zprávu druhého hráče.} \\
\\ 
\texttt{NÁVRH: \{'C08', 'C55', 'A67', 'B42', 'C43', 'C73', 'C10', 'B36', 'A18', 'C67', 'B40', 'C39', 'A21', 'C16'\}} \\
            }
        }
    }
    & \\ \\

    \theutterance \stepcounter{utterance}  
    & & & \multicolumn{4}{p{0.6\linewidth}}{
        \cellcolor[rgb]{0.9,0.9,0.9}{
            \makecell[{{p{\linewidth}}}]{
                \texttt{\tiny{[GM$\langle$P2]}}
                \texttt{SOUHLASÍM: \{'C08', 'C55', 'A67', 'B42', 'C43', 'C73', 'C10', 'B36', 'A18', 'C67', 'B40', 'C39', 'A21', 'C16'\}} \\
            }
        }
    }
     \\ \\

    \theutterance \stepcounter{utterance}  
    & & & \multicolumn{2}{p{0.3\linewidth}}{
        \cellcolor[rgb]{0.9,0.9,0.9}{
            \makecell[{{p{\linewidth}}}]{
                \texttt{\tiny{[GM$|$GM]}}
                \texttt{game successful} \\
            }
        }
    }
    & & \\ \\

    \theutterance \stepcounter{utterance}  
    & & & \multicolumn{2}{p{0.3\linewidth}}{
        \cellcolor[rgb]{0.9,0.9,0.9}{
            \makecell[{{p{\linewidth}}}]{
                \texttt{\tiny{[GM$|$GM]}}
                \texttt{end game} \\
            }
        }
    }
    & & \\ \\

\end{supertabular}
}

\end{document}
